\chapter{Modelle aus Mengenfamilien}
\label{b04guide:ch:models}

\section{Jedem Mengenglied seine Elemente zuordnen}

\begin{definition}[Mitgliedschaftsrelation einer Mengenfamilie]
Für eine Menge $\mathcal M$ von Mengen sei
\[
 \operatorname{Mem}(\mathcal M)
 =\{(X,x)\in\mathcal M\times\bigcup\mathcal M\mid x\in X\}.
\]
Dies ist die Mitgliedschaftsrelation $\MemRel(\mathcal M)$
aus Band~4.
\end{definition}

\begin{theorem}[Totalität genau bei nichtleeren Gliedern]
\label{b04guide:thm:membership}
Für jede Mengenfamilie $\mathcal M$ sind äquivalent:
\[
 \begin{aligned}
 &\TR\bigl(\operatorname{Mem}(\mathcal M);
                        \mathcal M,\bigcup\mathcal M\bigr),\\
 &\forall X\in\mathcal M\,(X\ne\varnothing),\\
 &\varnothing\notin\mathcal M.
 \end{aligned}
\]
Für jedes $X\in\mathcal M$ ist die Faser genau $X$.
Die leere Familie ist zugelassen. Die Mengenglieder müssen
nicht paarweise disjunkt sein.
\end{theorem}

\begin{definition}[Nichtleere Teilmengen als Ausgangsmenge]
Für eine Menge $B$ seien
\[
 \begin{aligned}
 \mathcal P_+(B)&=\mathcal P(B)\setminus\{\varnothing\},\\
 E_B&=\{(X,b)\in\mathcal P_+(B)\times B\mid b\in X\}.
 \end{aligned}
\]
\end{definition}

\begin{theorem}[Totale Mitgliedschaft auf der positiven Potenzmenge]
\label{b04guide:thm:positive}
Für jede Menge $B$ gilt
\[
 \TR(E_B;\mathcal P_+(B),B).
\]
Dabei ist $(E_B)_X=X$ für $X\in\mathcal P_+(B)$.
Für $B=\varnothing$ sind Ausgangsmenge und Relation beide leer.
\end{theorem}

\section{Elementen die sie enthaltenden Mengenglieder zuordnen}

\begin{definition}[Zugehörigkeit zu einer Mengenfamilie]
Für eine Menge $A$ und eine Mengenfamilie $\mathcal M$ sei
\[
 C_{A,\mathcal M}
  =\{(a,X)\in A\times\mathcal M\mid a\in X\}.
\]
\end{definition}

\begin{theorem}[Überdeckungskriterium]
\label{b04guide:thm:cover}
Für jede Menge $A$ und jede Mengenfamilie $\mathcal M$ gilt
\[
 \TR(C_{A,\mathcal M};A,\mathcal M)
 \quad\Longleftrightarrow\quad
 A\subseteq\bigcup\mathcal M.
\]
Insbesondere gilt stets
\[
 \TR\bigl(C_{\bigcup\mathcal M,\mathcal M};
                        \bigcup\mathcal M,\mathcal M\bigr).
\]
Leere oder sich überschneidende Glieder sind erlaubt.
\end{theorem}

\section{Aus nichtleeren Fasern zusammensetzen}

\begin{theorem}[Darstellung durch eine Menge von Faserangaben]
\label{b04guide:thm:fiberdata}
Seien $\mathcal H\subseteq A\times\mathcal P_+(B)$ und
\[
 \forall a\in A\,\exists X\in\mathcal P_+(B)\,
                                      ((a,X)\in\mathcal H).
\]
Dann ist
\[
 R_{\mathcal H}
  =\{(a,b)\in A\times B\mid
         \exists X\in\mathcal P_+(B)\,
                    ((a,X)\in\mathcal H\land b\in X)\}
\]
eine totale Relation von $A$ nach $B$.
Mehrere Mengenangaben zum selben $a$ sind zugelassen.
\end{theorem}

\begin{theorem}[Jede totale Relation besitzt eine solche Darstellung]
\label{b04guide:thm:representation}
Für $\TR(R;A,B)$ sei
\[
 \mathcal H_R=
 \{(a,X)\in A\times\mathcal P_+(B)\mid X=R_a\}.
\]
Dann gelten
\[
 \TR(\mathcal H_R;A,\mathcal P_+(B)),
 \qquad R_{\mathcal H_R}=R.
\]
Es werden die vollständigen nichtleeren Fasern erfaßt,
ohne aus ihnen einzelne Elemente auszuwählen.
\end{theorem}
